\chapter{Instalación paso a paso}
\label{cap:instalacion}

\begin{enpocas}
Qué hay que instalar y en qué orden: ROS~2 y Gazebo, el espacio de trabajo con nuestros paquetes, ArduPilot, el entorno de YOLO y las herramientas de CAD. Cada paso es un script que ya quedó en la carpeta personal.
\end{enpocas}

La instalación se hizo con cuatro scripts que quedaron en la carpeta personal
(\cod{~/instalar_ros2_gazebo.sh}, \cod{~/instalar_ardupilot.sh}, \cod{~/instalar_vision_yolo.sh}
y \cod{~/ardupilot_sitl.sh}). Esta sección explica qué hace cada uno y en qué orden se corren.

\section{ROS 2 Jazzy y Gazebo Harmonic}

\begin{lstlisting}[style=bash]
sudo bash ~/instalar_ros2_gazebo.sh
\end{lstlisting}

El script:
\begin{enumerate}
  \item agrega el repositorio \emph{universe} y el paquete \cod{ros2-apt-source} (que configura
        el repositorio oficial de ROS~2 para Ubuntu Noble);
  \item instala \cod{ros-jazzy-desktop}, \cod{ros-dev-tools}, el puente \cod{ros-jazzy-ros-gz}
        (trae Gazebo Harmonic), \cod{ros-jazzy-actuator-msgs}, \cod{ros-jazzy-xacro},
        PlotJuggler, el almacenamiento MCAP de \emph{rosbag2}, \cod{colcon} y numpy, scipy,
        matplotlib y pandas;
  \item inicializa \cod{rosdep}.
\end{enumerate}

Después, en cada terminal (o en \cod{~/.bashrc}):
\begin{lstlisting}[style=bash]
source /opt/ros/jazzy/setup.bash
\end{lstlisting}

Comprobación rápida: \cod{gz sim shapes.sdf} debe abrir Gazebo con unas figuras, y
\cod{ros2 topic list} debe responder.

\section{Espacio de trabajo}

\begin{lstlisting}[style=bash]
git clone https://github.com/jairoroke/Monitoreo-de-incendios-en-un-vehiculo-aereo-no-tripulado.git ~/ros2_ws
# Solo para controlador:=ardupilot (plugin de ArduPilot para Gazebo):
git clone https://github.com/ArduPilot/ardupilot_gazebo.git ~/ros2_ws/src/ardupilot_gazebo

cd ~/ros2_ws
export GZ_VERSION=harmonic
colcon build --symlink-install
source install/setup.bash
\end{lstlisting}

\cod{--symlink-install} enlaza los archivos de Python, los \emph{launch} y los YAML en lugar de
copiarlos: un cambio en \cod{src/} se ve sin volver a compilar (hay que recompilar solo si se
agrega un ejecutable nuevo en \cod{setup.py} o se cambia código C++). El espacio tiene dos
paquetes propios:

\begin{description}
  \item[\cod{cuadrirotor_sim}] (Python, \emph{ament\_python}): modelos, mundos, controladores,
        planificador, detector, \emph{launch}, scripts y pruebas.
  \item[\cod{cuadrirotor_gz_plugins}] (C++): dos sistemas de Gazebo, \cod{FuerzaViento} (fuerza
        externa de viento) y \cod{ComandosMotores} (para ArduPilot, sección~\ref{sec:plugin-ap}).
\end{description}

\section{ArduPilot SITL}
\label{sec:inst-ardupilot}

ArduCopter 4.7.2 se compila para SITL (\emph{software in the loop}) en \cod{~/ardupilot}:

\begin{lstlisting}[style=bash]
git clone --recurse-submodules https://github.com/ArduPilot/ardupilot.git ~/ardupilot
cd ~/ardupilot && git checkout Copter-4.7.2 && git submodule update --init --recursive
bash ~/instalar_ardupilot.sh          # como usuario normal, NO con sudo
./waf configure --board sitl && ./waf copter
\end{lstlisting}

\cod{instalar_ardupilot.sh} corre \cod{Tools/environment_install/install-prereqs-ubuntu.sh}
con dos variables que evitan conflictos con ROS:
\cod{DO_PYTHON_VENV_ENV=0} (no activar el entorno de Python de ArduPilot en cada terminal; chocaría
con el \cod{python3} de ROS) y \cod{DO_AP_STM_ENV=0} (no instalar el compilador para placas STM32).
Crea el entorno \cod{~/venv-ardupilot} con MAVProxy y pymavlink, e instala las bibliotecas de
GStreamer que necesita el plugin \cod{ardupilot_gazebo}. El binario queda en
\cod{~/ardupilot/build/sitl/bin/arducopter}.

En \cod{~/.bashrc} quedaron:
\begin{lstlisting}[style=bash]
export GZ_VERSION=harmonic
export GZ_SIM_SYSTEM_PLUGIN_PATH=$HOME/ros2_ws/install/ardupilot_gazebo/lib/ardupilot_gazebo:${GZ_SIM_SYSTEM_PLUGIN_PATH}
export GZ_SIM_RESOURCE_PATH=$HOME/ros2_ws/install/ardupilot_gazebo/share/ardupilot_gazebo/models:$HOME/ros2_ws/install/ardupilot_gazebo/share/ardupilot_gazebo/worlds:${GZ_SIM_RESOURCE_PATH}
export PATH=$HOME/ardupilot/Tools/autotest:$PATH
\end{lstlisting}

Prueba del SITL con el Iris de ejemplo (sin nuestra simulación):
\begin{lstlisting}[style=bash]
bash ~/ardupilot_sitl.sh              # iris_runway.sdf + sim_vehicle.py --console --map
\end{lstlisting}

Para la ventana que acomoda las ventanas del SITL (sección~\ref{sec:ventanas}) se agregó
\cod{python-xlib} al entorno de ArduPilot:
\begin{lstlisting}[style=bash]
~/venv-ardupilot/bin/pip install python-xlib
\end{lstlisting}

\section{Visión: YOLO26}

\begin{lstlisting}[style=bash]
bash ~/instalar_vision_yolo.sh        # descarga ~3 GB
\end{lstlisting}

Crea \cod{~/.venvs/yolo} con \cod{--system-site-packages} (el detector necesita ver los paquetes de
ROS~2) e instala \cod{ultralytics} con PyTorch y CUDA. Fija tres versiones:
\begin{itemize}
  \item \cod{numpy<2}: los mensajes de ROS~2 Jazzy están compilados con numpy 1.x;
  \item \cod{opencv-python<4.12}: las versiones nuevas exigen numpy 2;
  \item \cod{setuptools<80}: lo que acepta \cod{colcon}.
\end{itemize}
Al final comprueba que PyTorch vea la GPU y que \cod{rclpy} se importe desde el entorno.

\section{Herramientas de CAD y generación de mundos}

La conversión del CAD del Tarot y la generación del terreno usan bibliotecas que no vienen con
ROS. Van en un entorno aparte:
\begin{lstlisting}[style=bash]
python3 -m venv --without-pip ~/.venvs/cad
curl -sSL https://bootstrap.pypa.io/get-pip.py | ~/.venvs/cad/bin/python
~/.venvs/cad/bin/pip install cascadio trimesh numpy fast_simplification pillow scipy matplotlib
\end{lstlisting}

\begin{nota}
Se intentó primero convertir el STEP con \emph{gmsh}, pero tardaba más de 15 minutos;
\cod{cascadio} (teselador de OpenCASCADE) con \cod{trimesh} lo hace en segundos.
\end{nota}

\section{Comprobación de la instalación}

\begin{lstlisting}[style=bash]
cd ~/ros2_ws/src/cuadrirotor_sim && python3 -m pytest -q test/     # 31 pruebas
ros2 launch cuadrirotor_sim simulacion.launch.py                    # cuadrirrotor en hover
\end{lstlisting}
